\documentclass[11pt]{article}
\usepackage[a4paper,margin=1in]{geometry}
\usepackage{fontspec}
\usepackage[boldfont=false]{xeCJK}
\setCJKmainfont{FandolSong-Regular.otf}
\usepackage{amsmath}
\usepackage{amssymb}
\usepackage{array}
\usepackage{booktabs}
\usepackage{graphicx}
\usepackage{adjustbox}
\usepackage{enumitem}
\setlist[itemize]{topsep=2pt,partopsep=0pt,parsep=0pt,itemsep=2pt,leftmargin=1.4em}
\setlist[enumerate]{topsep=2pt,partopsep=0pt,parsep=0pt,itemsep=2pt,leftmargin=1.6em}
\usepackage{parskip}
\usepackage{fvextra}
\fvset{breaklines=true,breakanywhere=true,fontsize=\small}
\setlength{\parindent}{0pt}
\usepackage{longtable}
\setlength\LTpre{0pt}
\setlength\LTpost{\baselineskip}

\begin{document}

\begin{center}
  {\LARGE\bfseries Inheritance}\\[0.35em]
  {\large A Level Biology · Vocabulary}
\end{center}
\vspace{0.6em}

\begin{longtable}{@{}p{0.40\linewidth}p{0.24\linewidth}p{0.30\linewidth}@{}}
\toprule
\textbf{English} & \textbf{中文} & \textbf{Pinyin} \\
\midrule
\endhead
\bottomrule
\endlastfoot
diploid & 二倍体 & èr bèi tǐ \\
chromosome & 染色体 & rǎn sè tǐ \\
haploid & 单倍体 & dān bèi tǐ \\
homologous chromosomes & 同源染色体 & tóng yuán rǎn sè tǐ \\
gene & 基因 & jī yīn \\
gamete & 配子 & pèi zi \\
meiosis & 减数分裂 & jiǎn shù fēn liè \\
nuclear envelope & 核膜 & hé mó \\
spindle & 纺锤体 & fǎng chuí tǐ \\
crossing over & 交叉互换 & jiāo chā hù huàn \\
independent assortment & 自由组合 & zì yóu zǔ hé \\
fertilisation & 受精 & shòu jīng \\
protein & 蛋白质 & dàn bái zhì \\
locus & 基因座 & jī yīn zuò \\
allele & 等位基因 & děng wèi jī yīn \\
dominant & 显性 & xiǎn xìng \\
recessive & 隐性 & yǐn xìng \\
codominant & 共显性 & gòng xiǎn xìng \\
genotype & 基因型 & jī yīn xíng \\
phenotype & 表现型 & biǎo xiàn xíng \\
homozygous & 纯合 & chún hé \\
heterozygous & 杂合 & zá hé \\
test cross & 测交 & cè jiāo \\
Punnett square & 庞纳特方格 & páng nà tè fāng gé \\
monohybrid cross & 单基因杂交 & dān jī yīn zá jiāo \\
dihybrid cross & 双基因杂交 & shuāng jī yīn zá jiāo \\
multiple alleles & 复等位基因 & fù děng wèi jī yīn \\
sex linkage & 伴性遗传 & bàn xìng yí chuán \\
linkage & 连锁 & lián suǒ \\
autosome & 常染色体 & cháng rǎn sè tǐ \\
epistasis & 上位性 & shàng wèi xìng \\
chi-squared test & 卡方检验 & kǎ fāng jiǎn yàn \\
enzyme & 酶 & méi \\
albinism & 白化病 & bái huà bìng \\
haemoglobin & 血红蛋白 & xuè hóng dàn bái \\
sickle cell anaemia & 镰状细胞贫血 & lián zhuàng xì bāo pín xuè \\
haemophilia & 血友病 & xuè yǒu bìng \\
Huntington's disease & 亨廷顿病 & hēng tíng dùn bìng \\
gibberellin & 赤霉素 & chì méi sù \\
elongation & 伸长 & shēn cháng \\
structural gene & 结构基因 & jié gòu jī yīn \\
regulatory gene & 调节基因 & tiáo jié jī yīn \\
inducible enzyme & 可诱导酶 & kě yòu dǎo méi \\
repressible enzyme & 可抑制酶 & kě yì zhì méi \\
prokaryote & 原核生物 & yuán hé shēng wù \\
lac operon & 乳糖操纵子 & rǔ táng cāo zòng zi \\
repressor & 阻遏物 & zǔ è wù \\
transcription & 转录 & zhuǎn lù \\
transcription factor & 转录因子 & zhuǎn lù yīn zi \\
gene expression & 基因表达 & jī yīn biǎo dá \\
\end{longtable}

\end{document}
